% GDZ PPN236018647,native564–565before95;HKUST565–566.
\clearpage\SourceLeaf{564}\ArithmeticHead{535}{R}{1798 IUL.}
\NumberedItem{92}
\indent De lemniscata, elegantissima omnes exspectationes superantia acquisivimus et quidem per methodos, quae campum prorsus novum nobis aperiunt.
\par\noindent\hfill Gott[ingae, 1798] Iul.

\medskip\begingroup\small
\indent Die hier angezeigten Ergebnisse finden sich in der Scheda Aa, die im Juli 1798 begonnen worden ist. Unter der Überschrift: \textit{Scheda prima. De curva lemniscata} beginnt S. 3 der Scheda eine Zusammenfassung älterer Formeln, die im wesentlichen Werke III, S. 413—415 bis dahin, wo es heißt »Spätere Bemerkung« (diese stammt aus dem Handbuch 16, Bb, S. 72), abgedruckt ist. Das eigentlich Neue beginnt auf S. 6 der Scheda mit den Worten:
\par\begingroup\normalsize
\indent »Si valores sc [i. e. sinuscirculi], qui reddunt ipsum sl [i. e. sinumlemniscaticum]\(=0\) secundum regulas notas productum infinitum generare concipiuntur, nec non valores ipsius sc qui reddunt ipsum \(\mathrm{sl}=\infty\) .\ .\ .\ .\ .«
\par\endgroup
\noindent u.s.w., wie Werke III, S. 415, beginnend im art. [4.], gedruckt ist. \textsc{Gauss} knüpft also unmittelbar an die in der Nr. 91b des \textit{Tagebuchs} und in der Aufzeichnung oben S. 168, art. [4.], gegebenen Ansätze an, wonach \(\operatorname{sin\,lem}\varphi\) als Funktion von \(s=\sin\varphi\) betrachtet wird. Aus der Kenntnis der Null- und Unendlichkeitsstellen von \(\operatorname{sin\,lem}\varphi\) und \(\operatorname{cos\,lem}\varphi\) (siehe oben S. 153 ff., die artt. [2.], [5.], [6.]) ergeben sich dann unmittelbar die Werte von \(s\), die diesen Stellen entsprechen, in der Form von Exponentialgrößen, wie sie in den Ausdrücken für \(P\varphi,Q\varphi,p\varphi\) und \(q\varphi\) Werke III, S. 416 auftreten. \(P\varphi\) ist die oben S. 168, art. [3.] und S. 170, art. [6.] eingeführte Funktion; vergl. Werke III, S. 416, art. [5.], eine Aufzeichnung, die jedoch nicht aus der Scheda Aa, sondern von losen Zetteln aus späterer Zeit stammt.

Die jetzt gewonnene Darstellung von \(P\varphi,Q\varphi\) in der Form einfach unendlicher Produkte liefert auch die in der Tagebuchnotiz Nr. 63 (vom 29. März 1797) gemeinte \textit{demonstratio} des damals nur durch Induktion gefundenen Resultats (vergl. auch oben, S. 158 letzte Zeile)
\[
\theta=N\varpi=e^{\frac\pi2},
\]
\noindent sie brachte also wirklich, wie \textsc{Gauss} damals voraussagte, \textit{gravissima analyseos incrementa} mit sich. Daß \textsc{Gauss} diese Bestätigung sofort vorgenommen hat, zeigt eine in der Scheda Aa aufgezeichnete, aber Werke III nicht mit abgedruckte Tafel der Werte von \(e^{\frac12 k\pi}\) und \(e^{-\frac12 k\pi}\) für \(k=1,2,3,4,5,6\) und dann S. 8—13 der Scheda die Werke III, S. 431, beginnend mit »\textit{Ecce iam computum pro} \(e^{\frac12\pi}\)« bis S. 432 abgedruckte sehr genaue Berechnung von \(e^{\frac12\pi}\) und \(e^{-\frac12\pi}\). Dazwischen steht (noch auf S. 7 der Scheda) die Werke III, S. 417, art. [6.] wiedergegebene Partialbruchdarstellung und (auf S. 8 der Scheda) die den Schluß des art. [6.] bildende Entwicklung des sinus lemniscaticus, und zwar in der Form:
\[
\mathrm{sl}\,\varphi=\frac\pi\varpi\frac4{e^{\frac12\pi}+e^{-\frac12\pi}}\,\mathrm{sc}\,\varphi-\frac\pi\varpi\frac4{e^{\frac32\pi}-e^{-\frac32\pi}}\,\mathrm{sc}\,3\varphi\quad\text{etc.}
\]
\noindent Es ist dies also genau die in der Nr. 91b und oben S. 168 gegebene Entwicklung von \(\operatorname{sin\,lem}\varphi\) nach den Sinus der Vielfachen von \(\varphi\), nur erscheinen die Koeffizienten, die dort bloß numerisch angegeben waren, hier in einer das allgemeine Gesetz aufzeigenden analytischen Form. Aber auch hier hat \textsc{Gauss} auf beiden Seiten der Gleichung \(\varphi\) im Gradmaß gedacht (vergl. die Bemerkung bei der Nr. 91); beim Abdruck Werke III, S. 417 hat der Herausgeber \textsc{Schering} beiderseits auf Bogenmaß reduziert; die von ihm mit \(\psi\) bezeichnete
\par\endgroup

\clearpage\SourceLeaf{565}\ArithmeticHead{536}{L}{TAGEBUCH MIT ERLÄUTERUNGEN.}
\begingroup\small
\noindent Größe hat also den Wert
\[
\psi=\frac{\varphi^0}{180^0}.
\]

Mit denselben Änderungen gegen die Handschrift sind dann Werke III, S. 417, art. [7.] noch die auf den Seiten 14, 15 der Scheda stehenden Formeln für \(\log P,\ \log Q,\ \log p,\ \log q\) und \(\log\operatorname{sin\,lem}\varphi\) abgedruckt. Damit sind die zu der hier besprochenen Tagebuchnotiz Nr. 92 gehörigen Aufzeichnungen erschöpft. Vergl. die Bemerkungen zu den Nummern 94, 95, 98.
\par\noindent\hfill\textsc{Klein. Schlesinger.}
\par\endgroup

\NumberedItem{93}
\indent Solutio problematis ballistici.
\par\noindent\hfill Gott[ingae, 1798] Iul.

\medskip\begingroup\small
\indent Der Nachlaß enthält keine auf dieses Problem bezüglichen Aufzeichnungen. Dagegen hat \textsc{Gauss} in einem am 15. August 1842 an \textsc{Encke} geschriebenen Briefe einige Gedanken über die Art, wie ballistische \textls[120]{Versuche} anzustellen sind, ausgesprochen. Diesen Brief sowie die Anfrage \textsc{Enckes}, die ihm vorhergegangen war, findet man unter den Nachträgen zur Physik, Werke XI 1, S. 49 ff. abgedruckt.
\par\noindent\hfill\textsc{Schlesinger.}
\par\endgroup

\NumberedItem{94}
\indent Cometarum theoriam perfectiorem reddidi.
\par\noindent\hfill Gott[ingae, 1798] Iul.

\medskip\begingroup\small
\indent Auf die zur Theorie der lemniskatischen Funktionen gehörigen Entwicklungen der Scheda Aa, die bis zur S. 15 gehen und bei der Nr. 92 besprochen worden sind, folgt auf S. 16 die Überschrift: \textit{Scheda secunda. De motu cometarum.} Unter dieser Überschrift finden sich aber die Differentialgleichungen des Zweikörperproblems für die \textls[120]{Erde} und daran anschließende Entwicklungen, die sich bis zur S. 20 der Scheda erstrecken und kaum über die Ableitung des Flächensatzes und des Integrals der lebendigen Kraft hinausgehen; sie enthalten nichts, was sich im besonderen auf die Kometenbewegung bezieht. \textsc{Gauss}’ Untersuchungen über die Bahnbestimmung der Kometen fallen in eine spätere Zeit.
\par\noindent\hfill\textsc{Klein. Brendel.}
\par\endgroup
